% Figura: problema 4.8. Condensador esférico con dos dieléctricos hemisféricos.
\begin{figure}[t]
\centering
\begin{subfigure}[b]{0.5\textwidth}\centering
\begin{tikzpicture}[scale=1.0]
  \fill[dielec] (0,2.1) arc (90:270:2.1) -- cycle;
  \fill[dielecB] (0,-2.1) arc (-90:90:2.1) -- cycle;
  \draw[texto!70!fondo, line width=0.5pt] (0,-2.1) -- (0,2.1);
  \draw[texto, line width=1.1pt] (0,0) circle (2.1);
  \shade[ball color=conductorluz] (0,0) circle (0.85);
  \draw[texto, line width=0.6pt] (0,0) circle (0.85);
  \draw[gauss] (0,0) circle (1.5);
  \node[etiq] at (-0.9,1.62) {$\varepsilon_1$};
  \node[etiq] at (0.9,1.62) {$\varepsilon_2$};
  \draw[texto, line width=0.4pt, -{Stealth[length=1.6mm]}] (0,0) -- (40:0.85);
  \node[font=\small, text=black] at (0.12,0.12) [above left=-1pt] {$a$};
  \draw[acentoB, line width=0.5pt, -{Stealth[length=1.6mm]}] (0,0) -- (-25:1.5);
  \node[etiq, text=acentoB] at (-20:1.25) [below=1pt] {$r$};
  \draw[texto, line width=0.4pt, -{Stealth[length=1.6mm]}] (0,0) -- (-125:2.1);
  \node[etiq] at (-125:1.75) [left=1pt] {$b$};
  \node[etiq] at (-1.55,0.25) {$Q$};
  \draw[texto, line width=0.4pt, -{Stealth[length=1.5mm]}] (-1.38,0.18) -- (165:0.9);
  % tierra
  \draw[texto, line width=0.6pt] (2.1,0) -- (2.75,0) -- (2.75,-0.25);
  \pic at (2.75,-0.25) {tierra};
  \node[etiqt, anchor=west] at (2.82,0.15) {$\phi=0$};
\end{tikzpicture}
\caption{Geometría y superficie gaussiana}
\end{subfigure}\hfill
\begin{subfigure}[b]{0.47\textwidth}\centering
\begin{tikzpicture}[scale=0.42]
  \newcommand{\mini}[1]{\begin{scope}[shift={(0,#1)}]
     \fill[dielec] (0,2.1) arc (90:270:2.1) -- cycle;
     \fill[dielecB] (0,-2.1) arc (-90:90:2.1) -- cycle;
     \draw[texto, line width=0.6pt] (0,0) circle (2.1);
     \fill[fondo] (0,0) circle (0.95);\end{scope}}
  % (i) interfase: sin carga
  \mini{5.6}
  \draw[acentoC, line width=1.6pt] (0,0.95+5.6) -- (0,2.1+5.6) (0,-0.95+5.6) -- (0,-2.1+5.6);
  \draw[texto, line width=0.5pt] (0,5.6) circle (0.95);
  \node[etiq, anchor=west, align=left, font=\footnotesize] at (2.5,5.6) {interfase plana: $E_n=D_n=0$,\\ sin carga libre ni ligada};
  % (ii) carga total uniforme
  \mini{0}
  \draw[acentoA, line width=2.0pt] (0,0) circle (0.95);
  \node[etiq, anchor=west, align=left, font=\footnotesize] at (2.5,0) {carga total en $r=a$:\\ $\sigmaf+\sigmaP=\eo E(a)$, uniforme};
  % (iii) carga libre mayor en el medio de mayor permitividad
  \mini{-5.6}
  \draw[cargapos, line width=1.0pt] (0,0.95-5.6) arc (90:270:0.95);
  \draw[cargapos, line width=2.6pt] (0,-0.95-5.6) arc (-90:90:0.95);
  \node[etiq, anchor=west, align=left, font=\footnotesize] at (2.5,-5.6) {carga libre: $\sigma_{\mathrm{f},i}=\varepsilon_i E(a)$,\\ mayor donde $\varepsilon$ es mayor};
\end{tikzpicture}
\caption{Reparto de la carga en $r=a$ ($\varepsilon_2>\varepsilon_1$)}
\end{subfigure}
\caption{\Cref{prob:dos-hemisferios}. (a) Esfera interior de radio $a$ y carga $Q$; esfera exterior de radio $b$ a tierra; el hueco
está lleno por dos dieléctricos separados por un plano ecuatorial. (b) El campo es radial e igual en ambos
medios, pero la carga libre de la esfera interior no es uniforme: se concentra en el hemisferio del
dieléctrico de mayor permitividad, de modo que la carga total (libre más ligada) sí resulta uniforme.}
\label{fig:p48}
\end{figure}
